\section{Product covariances}
\label{sec:products}
The periodic construction includes products: if the physical mode at each site
is independent of the virtual modes, contraction leaves its covariance unchanged.
We make this explicit, allowing the physical covariance to vary from site to site.
We retain $N\geq3$, $M\geq1$ and the ordering of Section~\ref{sec:contraction}.

\begin{definition}\label{def:product-resource}
Let $\eta_i$ be a real one-mode covariance at site $i$. Take the local covariance
\[
 \Gamma_i=\diag(I_{2M},I_{2M},\eta_i)
\]
in grouped phase coordinates $(q_{A_i},p_{A_i},q_{B_i},p_{B_i},q_{C_i},p_{C_i})$,
then reorder it to the position-then-momentum coordinates used for periodic
assembly. Define the physical product covariance by
\[
 (\gamma_{\rm prod})_{(u,i),(v,j)}=
 \begin{cases}(\eta_i)_{uv},&i=j,\\0,&i\ne j,\end{cases}
 \qquad u,v\in\{q,p\}.
\]
\end{definition}
Each virtual mode is a vacuum and is independent of all other modes at its site.
The physical mode has covariance $\eta_i$.

\begin{theorem}\label{thm:periodic-product}
If every $\eta_i$ is admissible, then every $\Gamma_i$ is admissible and
\[
 W\gamma_{\rm loc}W^T=
 \begin{pmatrix}2I_{2NM}&0\\0&\gamma_{\rm prod}\end{pmatrix}.
\]
Consequently the periodic contraction equals $\gamma_{\rm prod}$ and is
admissible. In particular, one virtual mode per bond suffices for every $N\geq3$.
\end{theorem}
\begin{proof}
The independent-sum rule of Theorem~\ref{thm:sum} applies to the local resources.
In grouped phase coordinates its uncertainty matrix is
\[
 \Gamma_i-i\diag(\Omega_M,\Omega_M,\Omega_1)
 = (I_{2M}-i\Omega_M)\oplus(I_{2M}-i\Omega_M)
       \oplus(\eta_i-i\Omega_1)\succeq0.
\]
The two vacuum blocks and the prescribed physical block are positive semidefinite;
reordering the coordinates preserves the inequality.

Each eliminated quadrature is the signed sum of two independent virtual vacuum
coordinates. Its diagonal covariance entry is $1+1=2$. Distinct eliminated
coordinates have zero covariance: distinct virtual modes are independent, and
within each vacuum mode the position--momentum covariance vanishes. The virtual
coordinates are also independent of every physical coordinate. Thus, in the
notation of Definition~\ref{def:contraction},
\[
 E=2I_{2NM},\qquad F=0,\qquad G=\gamma_{\rm prod}.
\]
The Schur-complement formula therefore gives
\[
 \gamma_{\rm per}=G-F^TE^{-1}F=\gamma_{\rm prod}.
\]
Admissibility follows from Theorem~\ref{thm:periodic}.
\end{proof}

\begin{corollary}\label{cor:product-examples}
For $\eta_i=I_2$ at every site, the periodic contraction is $I_{2N}$.
For arbitrary $s_i>0$, the choices
$\eta_i=\diag(s_i^2,s_i^{-2})$ give the corresponding product of squeezed
one-mode covariances.
\end{corollary}
\begin{proof}
The vacuum covariance is admissible. The diagonal criterion of
Proposition~\ref{prop:diagonal} states
\[
 \diag(a,b)\text{ is admissible}\quad\Longleftrightarrow\quad a>0\ \text{and}\ ab\geq1.
\]
Apply it with $a=s_i^2$ and $b=s_i^{-2}$, whose product is one, and use
Theorem~\ref{thm:periodic-product}. In the vacuum case, the entries of
$\gamma_{\rm prod}$ are those of the identity.
\end{proof}
The one-mode Gaussians of Theorem~\ref{thm:covariance}, with $\hbar=1$, have
these squeezed covariances. The virtual vacua in this construction are unused;
adding them allows products to fit the same $M\geq1$ convention as the correlated
examples.